\section{Математическое описание}

\subsection{Экспертные системы}

Экспертная система~--- это программная система, использующая знания из некоторой предметной области для решения задач, которые обычно требуют участия специалиста. В данной работе экспертная система имитирует консультацию специалиста по подготовке макета бейджа к печати: пользователь отвечает на вопросы, а система по этим ответам выбирает подходящие параметры шрифтовой гарнитуры.

\subsection{Структура экспертной системы}

Типовая экспертная система включает базу знаний, рабочую память, механизм логического вывода и пользовательский интерфейс. База знаний хранит правила предметной области, рабочая память содержит сведения, полученные в ходе консультации, а механизм вывода выбирает применимые правила и формирует итоговое заключение.

\subsection{Механизм логического вывода}

Механизм логического вывода выполняет переход от исходных сведений к итоговому заключению. В данной работе используется прямой логический вывод: система последовательно сопоставляет накопленные ответы с условиями правил и при достижении листа дерева формирует набор рекомендаций.

\subsection{Продукционная модель представления знаний}

Продукционная модель представления знаний описывает знания в виде правил:
\[
    \text{ЕСЛИ } <\text{условие}> \text{, ТО } <\text{действие}>.
\]
Условие правила называется антецедентом, а действие~--- консеквентом. Для данной задачи антецедент задаёт сочетание признаков макета и условий печати, а консеквент содержит переход к следующему вопросу или итоговую рекомендацию.

В общем виде продукцию можно представить как:
\[
    N = \langle A, U, C, I, R \rangle,
\]
где $N$~--- имя продукции, $A$~--- область её применения, $U$~--- условие применимости, $C$~--- ядро продукции, $I$~--- постусловие, а $R$~--- пояснение или служебная информация.

Достоинством продукционной модели является простота добавления и чтения отдельных правил. Каждое правило можно рассматривать независимо, поэтому небольшая экспертная система остаётся понятной и легко расширяется. Недостатки проявляются при росте числа правил: становится сложнее контролировать полноту базы знаний, пересечения условий и общий порядок вывода.

\subsection{Бинарное дерево решений}
\label{sec:decision-tree-model}

Бинарное дерево решений~--- это ориентированный ациклический граф, в котором каждый внутренний узел соответствует вопросу с двумя вариантами ответа, а каждый лист соответствует конечному решению. В данной работе дерево решений используется как наглядная форма организации консультации. Путь от корня к листу задаёт последовательность вопросов, а лист содержит набор рекомендаций по параметрам шрифта.

Внутренними критериями дерева являются:
\begin{itemize}
    \item наличие дополнительного текста помимо имени и логотипа;
    \item наличие кириллицы;
    \item необходимость чтения с расстояния более одного метра;
    \item тёмный фон бейджа;
    \item наличие ламинации;
    \item плотность бумаги выше 250 г/м².
\end{itemize}

Листья дерева соответствуют разным сочетаниям этих критериев и содержат итоговые рекомендации.
